\begin{alist}
\item $S(0) = 0$ (το κινητό βρίσκεται στην αφετηρία) και $S(2) = 2 \cdot 8 - 6 \cdot 4 + 20 = 12 \text{ μέτρα}$.

\item $S(t) = 30 \Leftrightarrow 2t^3 - 6t^2 + 10t = 30 \Leftrightarrow t^3 - 3t^2 + 5t - 15 = 0 \quad (1)$.
Πιθανές ακέραιες ρίζες της (1) είναι οι $\pm 1, \pm 3, \pm 5, \pm 15$. Επειδή ο χρόνος είναι μη αρνητικό φυσικό μέγεθος, πιθανές ρίζες είναι οι $1, 3, 5, 15$.
Η τιμή $t = 3$ επαληθεύει την εξίσωση και με τη βοήθεια του σχήματος Horner, έχουμε:

\begin{center}
\begin{tabular}{|c|c|c|c|c|}
\hline
1 & -3 & 5 & -15 & \cellcolor{green!50} 3 \\ \hline
\cellcolor{gray!20} & 3 & 0 & 15 & \multicolumn{1}{c}{} \\ \cline{1-4}
\cellcolor{cyan!50} 1 & \cellcolor{cyan!50} 0 & \cellcolor{cyan!50} 5 & \cellcolor{cyan!50} 0 & \multicolumn{1}{c}{} \\ \cline{1-4}
\end{tabular}
\end{center}

Άρα η $(1) \Leftrightarrow (t - 3)(t^2 + 5) = 0 \Leftrightarrow t = 3$.
Επομένως το κινητό θα χρειαστεί 3 δευτερόλεπτα για να διανύσει απόσταση 30 μέτρων.

\item Πραγματικά, έχουμε:
$S(t) = 2t^3 - 6t^2 + 10t = 2t(t^2 - 3t + 5) \ge 0$,
διότι $t \ge 0$ (εκφράζει το χρόνο)
και $t^2 - 3t + 5 > 0$ (το τριώνυμο έχει διακρίνουσα $\Delta = -11 < 0$, επομένως είναι πάντα ομόσημο του συντελεστή του $t^2$).

\item Με βάση το φυσικό πλαίσιο του προβλήματος, η συνάρτηση $S(t)$ πρέπει να είναι μη αρνητική και σε κανένα χρονικό διάστημα γνήσια φθίνουσα. Επομένως, είναι:

Η (I) είναι μεν μη αρνητική, αλλά δε διατηρεί το ίδιο είδος μονοτονίας.

Η (II) παίρνει και αρνητικές τιμές.

Η (III) είναι μη αρνητική και γνησίως αύξουσα παντού, ως εκ τούτου αποτελεί την ενδεδειγμένη απάντηση.
\end{alist}
